% ------------------------------------------------------------------------
% file `c2-exercise-3-solution-body.tex'
%   in folder `xsim/'
%
%     solution of type `exercise' with id `3'
%
% generated by the `solution' environment of the
%   `xsim' package v0.21 (2022/02/12)
% from source `c2' on 2023/06/13 on line 88
% ------------------------------------------------------------------------
    Évaluations :
    \begin{itemize}
        \item Le raisonnement est correcte et rigoureux\hfill 0,5
        \item La réponse est correcte \hfill 0,5
    \end{itemize}
\begin{tasks}
    \task
    Comme le triangle $ABC$ est rectangle en $B$,
    \begin{itemize}
        \item
        $
        \begin{aligned}[t]
            |AC|^2 &= |BC|^2 + |AB|^2\\
            |AC| &= \sqrt{|BC|^2 + |AB|^2}\\
            |AC| &= \sqrt{\left(\sqrt{32}\right)^2 + \left(\sqrt{17}\right)^2}\\
            |AC| &= \sqrt{49}\\
            |AC| &= 7
        \end{aligned}
        $
        \item
        $
        \begin{aligned}[t]
            |\widehat{C}| &= \arctan\left(\dfrac{\sqrt{17}}{\sqrt{32}}\right)\\
            |\widehat{C}| &= \ang{36,1}
        \end{aligned}
        $
    \end{itemize}

    \task
    Comme le triangle $ABC$ est rectangle en $C$,
    \begin{itemize}
        \item
        $
        \begin{aligned}[t]
            |AC| &= 4 \cos\left(\ang{25}\right)\\
            |AC| &= 3,6
        \end{aligned}
        $
        \item
        $
        \begin{aligned}[t]
            |BC| &= 4\sin\left(\ang{25}\right)\\
            |BC| &= 1,7
        \end{aligned}
        $
    \end{itemize}

    \task
    Comme $DE \sslash BC$,
    \begin{itemize}
        \item
        $
        \begin{aligned}[t]
            \dfrac{|AD|}{|AE|} &= \dfrac{|AB|}{|AC|}\\
            \dfrac{|AD|}{24} &= \dfrac{7}{10}\\
            |AD| &= 24 \cdot \dfrac{7}{10}\\
            |AD| &= \dfrac{84}{5} = 16,8
        \end{aligned}
        $

        \item
        $
        \begin{aligned}[t]
          \dfrac{|AE|}{|AC|} &= \dfrac{|DE|}{|BC|}\\
          \dfrac{24}{10} &= \dfrac{12}{|BC|}\\
          |BC| &= 12 \cdot \dfrac{10}{24}\\
          |BC| &= 5
        \end{aligned}
        $
    \end{itemize}

    \task
    \begin{itemize}
        \item
        $O$ est le centre du cercle, $A$ et $B$ sont deux points sur le cercle. Donc, $[OB]$ et $[OA]$ sont deux rayons. Par conséquent, le triangle $AOB$ est isocèle en $O$. On en déduit que $|\widehat{O}| = 180-2\cdot 40 = \ang{100}$.

        \item
        $\widehat{C}$ est un angle inscrit (car $C$ appartient au cercle) et $\widehat{O}$ est un angle au centre. Ils interceptent l'arc $\wideparen{AB}$. Donc, $|\widehat{C}| = \dfrac{|\widehat{O}|}{2}=\ang{50}$.
    \end{itemize}
\end{tasks}
